\section{Introduction}
\label{sec:introduction}

For a smooth projective complex variety \(X\), the anticanonical line
bundle is \(-K_X=\det T_X\). A smooth Hermitian metric on this line is
semipositive when its Chern curvature is a nonnegative real \((1,1)\)-form.
Such a metric makes the numerical class nef. The nonvanishing question
asks whether \(H^0(X,-mK_X)\ne0\) for some integer \(m>0\), or
equivalently whether some positive multiple of \(-K_X\) is linearly
equivalent to an effective divisor.

We prove the following nonvanishing theorem.

\begin{theorem}[Anticanonical nonvanishing]\label{thm:headline}
Let \(X\) be a smooth connected projective variety over \(\C\).
If \(-K_X\) admits a smooth Hermitian metric of semipositive curvature,
then \(H^0(X,-mK_X)\ne0\) for some integer \(m>0\).
\end{theorem}

This is the positive-multiple formulation associated with Yau's
anticanonical nonvanishing question \cite[Problem~75, p.~395]{Yau1994}.
The multiple matters even when the anticanonical bundle is not torsion.
For a complex Enriques surface \(E\), \(K_E\) is nontrivial of order two
and has no section \cite[Section~3]{Dolgachev2016}.
Write \(p_E\) and \(p_{\mathbb P^1}\) for the two projections from
\(E\times\mathbb P^1\). Its anticanonical line is
\[
 p_E^*(-K_E)\otimes p_{\mathbb P^1}^*\OO_{\mathbb P^1}(2).
\]
The flat metric on the first factor and Fubini--Study metric on the
second make it smoothly semipositive. The product formula for sections
gives no first-power section, whereas the square has sections.
Restriction to a \(\mathbb P^1\) fiber shows that this anticanonical
line is not torsion.

\subsection{Geometry, monodromy, and earlier methods}

The structure theory explains both the strength of the metric assumption
and the difficulty that remains. Yau's prescribed-Ricci theorem
\cite{Yau1978} realizes the given smooth semipositive anticanonical
curvature as the Ricci form of a K\"ahler metric. The structure results
of Demailly--Peternell--Schneider \cite[Structure Theorem, p.~218]{DPS1996} and
Campana--Demailly--Peternell \cite[Theorem~1.4]{CDP2015} describe the
universal cover by its flat, compact Ricci-flat, and remaining compact
factors. The remaining factor is projective and has no holomorphic forms
of positive degree. After passage to a finite-index deck subgroup, the
actions that remain on this factor lie in a compact torus.

This last action cannot be ignored. A section on the compact factor
descends only if it is compatible with the deck transformations.
The flat and Ricci-flat factors supply canonical frames; on the remaining
factor we construct an anticanonical section fixed by the torus, for the
natural action on that factor's anticanonical bundle. Its product with those frames descends
to a finite \'etale cover, and its finite \'etale norm then
gives a section on \(X\). This reasoning uses the action on the universal
cover, not a global product decomposition of a finite cover.

In the broader nef setting,
Lazi\'c--Matsumura--Peternell--Tsakanikas--Xie proved numerical effectivity
for projective log canonical threefold pairs with rational boundary
and nef log-anticanonical divisor, under
\(\mathbb Q\)-factoriality or rational singularities
\cite[Theorem~A]{LMPTX2023}. Numerical effectivity means that the divisor
class is numerically equivalent to an effective rational divisor;
the existence of a section requires rational \emph{linear} effectivity.
For projective klt \(\mathbb Q\)-pairs with nef log-anticanonical
divisor, M\"uller proved actual nonvanishing in dimension three,
and in arbitrary dimension when the log-anticanonical divisor of a
general fiber of the maximal rationally connected quotient is semiample
\cite[Theorem~A and Corollary~B]{Muller2025}.
His equivariant theorem constructs naturally invariant plurisections
of a semiample log-anticanonical divisor for boundary-preserving actions
of commutative linear algebraic groups on projective sub-log-canonical
\(\mathbb Q\)-pairs
\cite[Theorem~C]{Muller2025}. It is a direct methodological predecessor
for addressing residual monodromy. Here the finite-volume argument
constructs the needed naturally invariant section from smooth
semipositivity, without assuming that fiber semiampleness.

In dimension two, Chen--Filip--Sun--Tosatti--Zhang have classified
compact K\"ahler surfaces with Hermitian-semipositive anticanonical bundle
\cite[Theorem~1.3]{CFSTZ2026}.
Their classification gives complementary low-dimensional geometry;
the present argument develops a section-producing mechanism in arbitrary
dimension.\footnote{The authors report that ChatGPT~5.6~Sol suggested
a preliminary version of their Lemma~5.2, which they reformulated and
developed. This records their declared assistance, not an attribution
of authorship.}

The ordinary part of the proof builds on two further developments.
Generic positivity and foliation methods control cotangent tensors:
the relevant inputs come from movable-cone duality
\cite{BDPP2013}, movable slope theory \cite{GKP2016}, and
Campana--P\u aun's positive-slope algebraicity criterion
\cite[Theorem~1.1]{CampanaPaun2019}.
The slope argument follows the foliation strategy used by
Ou \cite[Theorem~1.4]{Ou2023} and by
Lazi\'c--Matsumura--Peternell--Tsakanikas--Xie
\cite[Theorem~4.1]{LMPTX2023}.
Here the relative Bergman metric supplies ramification-corrected
positivity when an effective boundary and a finite adjoint measure
replace nef anticanonical data.
Hard Lefschetz with multiplier ideals \cite[Theorem~2.1.1]{DPS2001}
then produces twisted forms, whose determinants are the effective
divisors used to select a smaller base.

The transfer in Section~\ref{sec:transfer} uses the relative Bergman metric of
Berndtsson--P\u aun \cite[Theorem~0.1]{BerndtssonPaun2008}.
The ordinary proof uses its singular-coefficient form. The equivariant
proof instead uses Berndtsson's smooth direct-image positivity
\cite[Theorem~1.2]{Berndtsson2009}: averaging is an orthogonal holomorphic
projection onto the invariant line. These are distinct applications of
positivity, and we state their hypotheses separately.

\subsection{A finite-volume theorem and the proof strategy}
\label{sec:overview}

The inductive statement involves two line bundles. We use additive
notation: $J+L$ means $J\otimes L$, and $mL$ means $L^{\otimes m}$.
For an effective integral divisor $D$, write $s_D$ for the canonical
section of $\OO_Z(D)$. If $J=-K_Z+D$, this section is a $J$-valued
holomorphic top form. A metric $h_J$ therefore pairs $s_D$ with its
conjugate to give a measure, denoted $|s_D|_{h_J}^2$.

A singular metric has local squared frame norm $e^{-\phi}$. It is
semipositive when $\phi$ is plurisubharmonic; locally bounded weights
means bounds on both sides.

\begin{theorem}[Finite volume and sections with prescribed poles]
\label{thm:finite-volume}
Let $Z$ be a smooth connected projective complex variety of dimension
$d$ with $H^0(Z,\Omega_Z^q)=0$ for every $1\le q\le d$.
Let $D\ge0$ be an integral divisor. Suppose that $J=-K_Z+D$ has a
semipositive singular Hermitian metric $h_J$ satisfying
\[
                 \int_Z |s_D|_{h_J}^2<\infty.
\]
If a line bundle $L$ on $Z$ has a semipositive singular Hermitian metric
with locally bounded weights, then there exist integers $a>0$ and
$b\ge0$ such that
\[
                       H^0(Z,aL+bD)\ne0.
\]
\end{theorem}

This generality is needed because a fibration need not transfer the
line whose section is sought to the anticanonical bundle of its base.
Its fiber integrals instead produce a new target line on the base and
a separate corrected anticanonical line carrying the volume metric.
The volume metric $h_J$ may be unbounded, provided the indicated
measure is integrable. The metric on $L$ has the separate boundedness
assumption. The conclusion says that a positive power of $L$
has a rational section with poles only on $D$; the pole order $b$ is not
uniform. Rationally connected varieties satisfy the no-forms hypothesis,
but the theorem also applies more generally. For example, an Enriques
surface has no positive-degree holomorphic forms and a torsion canonical
bundle, although it is not rationally connected.

Theorem~\ref{thm:finite-volume} is proved by induction on dimension.
We first describe how its components fit together.
Assume that its conclusion fails. Then $L+D$ is pseudoeffective but not
big. A supporting movable curve class $\alpha\ne0$ satisfies
$L\cdot\alpha=D\cdot\alpha=0$. Finite adjoint volume bounds the slopes
of cotangent tensors against $\alpha$; in particular their determinant
lines have nonpositive degree. This is the content of
Section~\ref{sec:slope}.

The absence of holomorphic forms gives $\chi(\OO_Z)=1$. Riemann--Roch
and hard Lefschetz then produce sections
\[
 s_i\in H^0(Z,M+m_iL),\qquad m_i\longrightarrow\infty,
\]
where $M$ is one fixed cotangent determinant with $M\cdot\alpha\le0$.
Consider all nonempty complete systems $|uL+vM+wD|$ with nonnegative
integral coefficients. None can define a generically finite map: that
would make its divisor big, contradicting its nonpositive pairing with
$\alpha$. Choose a system with maximal image dimension and take the
relative algebraic closure of its image field in $\C(Z)$. This gives a
rational fibration onto a strictly smaller smooth projective base $S$.
Products of sections force every ratio in the indicated systems to be
a function on $S$.

This control of section ratios has two consequences. First, ratios among the
$s_i$ give a rational section $\rho$ of some $kL$, $k>0$, with no poles
on divisors dominating $S$. Second, on a smooth resolution $Y$ with maps $\pi:Y\to Z$ and
$f:Y\to S$, put $K_{Y/S}=K_Y-f^*K_S$. The relative adjoint spaces for
$K_{Y/S}+\pi^*(J+jkL)$ have dimension one for every $j\ge0$.
They are generated over the base field by the pulled-back form $s_D$
with its birational Jacobian,
divided by a base canonical frame and multiplied by $\rho^j$.
We prove transfer first under these explicit hypotheses;
Section~\ref{sec:induction} then constructs its inputs and closes the
induction, including the extension and common denominator arguments.

Section~\ref{sec:transfer} explains why these algebraic facts are the
right ones. On a normal equidimensional model, minima of vertical
divisor orders turn $\rho$ into a regular section $e$ after subtracting a
base line $A$. On a smooth resolution, integrate its powers against
the pulled-back adjoint measure. The zeroth integral gives a metric on
$-K_S+D_S$ with finite adjoint volume. Normalized high moments give a
locally bounded semipositive metric on $A$. Normalization leaves $e$
nonzero on some component above every base divisor. When integrating
$e^j$, the boundary estimate therefore has a fixed order, rather than
one growing with $j$; dividing its logarithm by $j$ removes that
boundary contribution. The correction $D_S$ is
supported on base primes for which every component upstairs is exceptional over
$Z$ or lies above $D$. Thus the induction hypothesis on $S$ yields a
section that lifts with no forbidden pole on $Z$.

To obtain the torus-invariant section needed by the geometric reduction,
Section~\ref{sec:torus} first proves the smooth invariant transfer
statement, Proposition~\ref{prop:invariant-transfer}, and then applies
it to a rational orbit space.
Fundamental vector fields give the relative section, and a dense orbit
gives invariant adjoint rank one in every degree. Smooth direct-image
positivity applies to the invariant line because compact averaging is
an orthogonal holomorphic projection. The ordinary theorem on the
quotient base then produces a naturally invariant anticanonical section.
The cover and norm results of Section~\ref{sec:structure} finish the
proof of Theorem~\ref{thm:headline}.

The proofs below establish the model preparation, slope bound, and
rank and transfer arguments using the cited geometric and analytic results.
Section~\ref{sec:models} introduces the two birational models and the
metric conventions used throughout.
